\documentclass[12pt,letterpaper]{article}
\usepackage[utf8]{inputenc}
\usepackage[spanish]{babel}
\usepackage{amsmath,amssymb}
\usepackage{graphicx}
\usepackage{xcolor}
\usepackage{tikz}
\usepackage{booktabs}
\usepackage{array}
\usepackage{multirow}
\usepackage{fancyhdr}
\usepackage{geometry}
\usepackage{hyperref}
\usepackage{enumitem}
\usepackage{tcolorbox}

\geometry{top=2.5cm, bottom=2.5cm, left=2.5cm, right=2.5cm}

% Colores corporativos
\definecolor{primary}{RGB}{99, 102, 241}
\definecolor{secondary}{RGB}{16, 185, 129}
\definecolor{accent}{RGB}{245, 158, 11}
\definecolor{darkbg}{RGB}{15, 23, 42}
\definecolor{lighttext}{RGB}{148, 163, 184}

% Configuracion de hyperref
\hypersetup{
    colorlinks=true,
    linkcolor=primary,
    urlcolor=primary
}

% Header y footer
\pagestyle{fancy}
\fancyhf{}
\fancyhead[L]{\textcolor{primary}{\textbf{ElectroAI}}}
\fancyhead[R]{\textcolor{lighttext}{An\'alisis de Despliegue}}
\fancyfoot[C]{\thepage}
\renewcommand{\headrulewidth}{0.5pt}

\title{
    \vspace{-1cm}
    \textcolor{primary}{\rule{\linewidth}{2pt}}\\[0.5cm]
    {\Huge\textbf{An\'alisis de Factibilidad}}\\[0.3cm]
    {\LARGE Despliegue de ElectroAI en Infraestructura Universitaria}\\[0.5cm]
    \textcolor{primary}{\rule{\linewidth}{2pt}}
}
\author{
    \textbf{ElectroAI - Asistente de Electromagnetismo}\\
    Universidad del B\'io-B\'io
}
\date{\today}

\begin{document}

\maketitle
\thispagestyle{empty}

\vspace{1cm}

\begin{tcolorbox}[colback=primary!10, colframe=primary, title=\textbf{Resumen Ejecutivo}]
Este documento analiza la factibilidad de desplegar ElectroAI utilizando infraestructura local (rack de GPUs) en lugar de APIs comerciales. La conclusi\'on principal es que \textbf{el despliegue local es factible, m\'as econ\'omico y proporciona mayor control} sobre el servicio educativo.
\end{tcolorbox}

\newpage
\tableofcontents
\newpage

%========================================
\section{Contexto y Problema}
%========================================

\subsection{Situaci\'on Actual}

ElectroAI es una aplicaci\'on de asistencia educativa para electromagnetismo que utiliza:
\begin{itemize}
    \item \textbf{Claude Sonnet 4} (Anthropic) como modelo de lenguaje
    \item \textbf{ChromaDB} para almacenamiento vectorial
    \item \textbf{Streamlit} como interfaz web
    \item \textbf{Sistema RAG} para recuperaci\'on de problemas resueltos
\end{itemize}

\subsection{Problema Identificado}

El uso de la API de Anthropic presenta las siguientes limitaciones:
\begin{enumerate}
    \item \textbf{Costo por consulta}: Cada estudiante requiere cr\'editos API
    \item \textbf{Dependencia externa}: El servicio depende de disponibilidad de Anthropic
    \item \textbf{Privacidad}: Los datos de consultas viajan a servidores externos
    \item \textbf{Escalabilidad limitada}: Costos crecen linealmente con el uso
\end{enumerate}

\subsection{Recurso Disponible}

La universidad cuenta con un \textbf{rack de tarjetas gr\'aficas de \'ultima generaci\'on} disponible para aplicaciones de inteligencia artificial.

%========================================
\section{An\'alisis T\'ecnico}
%========================================

\subsection{Limitaci\'on de Claude}

\begin{tcolorbox}[colback=red!10, colframe=red!70!black, title=\textbf{Importante}]
\textbf{Claude NO puede ejecutarse localmente.} Es un modelo propietario cerrado de Anthropic, disponible \'unicamente v\'ia API de pago.
\end{tcolorbox}

\subsection{Alternativa: Modelos Open Source}

Existen modelos de c\'odigo abierto con calidad comparable que \textbf{s\'i pueden ejecutarse en hardware local}:

\begin{table}[h]
\centering
\caption{Modelos Open Source Recomendados}
\begin{tabular}{@{}lccl@{}}
\toprule
\textbf{Modelo} & \textbf{Par\'ametros} & \textbf{VRAM} & \textbf{Calidad} \\
\midrule
Llama 3.1 70B & 70B & 40 GB & Excelente \\
Llama 3.1 8B & 8B & 8 GB & Muy buena \\
Mistral 7B & 7B & 6 GB & Buena \\
Qwen 2.5 72B & 72B & 45 GB & Excelente \\
DeepSeek-V2 & 236B (MoE) & 80 GB & Superior \\
\bottomrule
\end{tabular}
\end{table}

\subsection{Servidores de Inferencia}

Para servir modelos localmente se pueden usar:

\begin{table}[h]
\centering
\caption{Opciones de Servidor de Inferencia}
\begin{tabular}{@{}llll@{}}
\toprule
\textbf{Software} & \textbf{Facilidad} & \textbf{Rendimiento} & \textbf{Uso} \\
\midrule
Ollama & F\'acil & Bueno & Desarrollo/Producci\'on peque\~na \\
vLLM & Media & Excelente & Producci\'on \\
TGI (HuggingFace) & Media & Muy bueno & Producci\'on \\
\bottomrule
\end{tabular}
\end{table}

%========================================
\section{Arquitectura Propuesta}
%========================================

\begin{center}
\begin{tikzpicture}[
    node distance=1.5cm,
    box/.style={rectangle, draw=primary, fill=primary!10, rounded corners, minimum width=4cm, minimum height=1cm, align=center},
    smallbox/.style={rectangle, draw=secondary, fill=secondary!10, rounded corners, minimum width=2.5cm, minimum height=0.8cm, align=center, font=\small},
    arrow/.style={->, thick, primary}
]

% Servidor GPU
\node[box, minimum width=10cm, minimum height=2.5cm] (gpu) at (0,4) {};
\node[above] at (gpu.north) {\textbf{Servidor GPU (Rack Universidad)}};
\node[smallbox] at (-2,4.3) {Llama 3.1 70B};
\node[smallbox] at (2,4.3) {Ollama/vLLM};
\node[font=\small, lighttext] at (0,3.5) {API REST: puerto 11434};

% Aplicacion
\node[box, minimum width=10cm, minimum height=2cm] (app) at (0,1.5) {};
\node[above] at (app.north) {\textbf{ElectroAI Application}};
\node[smallbox, fill=accent!20, draw=accent] at (-3,1.5) {Streamlit};
\node[smallbox, fill=accent!20, draw=accent] at (0,1.5) {RAG System};
\node[smallbox, fill=accent!20, draw=accent] at (3,1.5) {ChromaDB};

% Proxy
\node[box, minimum width=6cm] (proxy) at (0,-0.5) {Nginx + SSL + Auth LDAP};

% Estudiantes
\node[smallbox] at (-4,-2.5) {Estudiante 1};
\node[smallbox] at (0,-2.5) {Estudiante 2};
\node[smallbox] at (4,-2.5) {Estudiante N};

% Flechas
\draw[arrow] (gpu) -- (app);
\draw[arrow] (app) -- (proxy);
\draw[arrow] (proxy) -- (-4,-1.8);
\draw[arrow] (proxy) -- (0,-1.8);
\draw[arrow] (proxy) -- (4,-1.8);

\end{tikzpicture}
\end{center}

\subsection{Componentes del Sistema}

\begin{enumerate}
    \item \textbf{Capa de Inferencia}: Servidor GPU ejecutando modelo Llama 3.1
    \item \textbf{Capa de Aplicaci\'on}: ElectroAI con RAG y ChromaDB
    \item \textbf{Capa de Acceso}: Nginx con SSL y autenticaci\'on LDAP universitaria
    \item \textbf{Clientes}: Navegadores web (m\'ovil y escritorio)
\end{enumerate}

%========================================
\section{An\'alisis de Costos}
%========================================

\subsection{Escenario de Uso}

\begin{tcolorbox}[colback=secondary!10, colframe=secondary]
\textbf{Par\'ametros del escenario:}
\begin{itemize}[noitemsep]
    \item 500 estudiantes activos
    \item 50 consultas promedio por estudiante/mes
    \item 12 meses de operaci\'on
    \item Total: \textbf{300,000 consultas/a\~no}
\end{itemize}
\end{tcolorbox}

\subsection{Opci\'on 1: API de Anthropic (Claude)}

\begin{table}[h]
\centering
\caption{Costos API Anthropic - Claude Sonnet 4}
\begin{tabular}{@{}lr@{}}
\toprule
\textbf{Concepto} & \textbf{Costo} \\
\midrule
Precio Input & \$3.00 / 1M tokens \\
Precio Output & \$15.00 / 1M tokens \\
\midrule
Tokens promedio por consulta (input) & 2,000 tokens \\
Tokens promedio por consulta (output) & 1,000 tokens \\
\midrule
Costo por consulta (input) & \$0.006 \\
Costo por consulta (output) & \$0.015 \\
\textbf{Costo total por consulta} & \textbf{\$0.021} \\
\midrule
\textbf{Costo mensual (25,000 consultas)} & \textbf{\$525} \\
\textbf{Costo anual (300,000 consultas)} & \textbf{\$6,300} \\
\bottomrule
\end{tabular}
\end{table}

\subsection{Opci\'on 2: Hardware Local}

\begin{table}[h]
\centering
\caption{Costos Despliegue Local}
\begin{tabular}{@{}lr@{}}
\toprule
\textbf{Concepto} & \textbf{Costo Anual} \\
\midrule
Hardware (ya disponible) & \$0 \\
Electricidad (estimado 500W promedio, 8h/d\'ia) & \$350 \\
Mantenimiento y soporte & \$500 \\
Licencias software (todo open source) & \$0 \\
\midrule
\textbf{Total anual} & \textbf{\$850} \\
\bottomrule
\end{tabular}
\end{table}

\subsection{Comparativa y Ahorro}

\begin{table}[h]
\centering
\caption{Comparativa de Costos (5 a\~nos)}
\begin{tabular}{@{}lrrr@{}}
\toprule
\textbf{A\~no} & \textbf{API Anthropic} & \textbf{Local} & \textbf{Ahorro} \\
\midrule
A\~no 1 & \$6,300 & \$850 & \$5,450 \\
A\~no 2 & \$6,300 & \$850 & \$5,450 \\
A\~no 3 & \$6,300 & \$850 & \$5,450 \\
A\~no 4 & \$6,300 & \$850 & \$5,450 \\
A\~no 5 & \$6,300 & \$850 & \$5,450 \\
\midrule
\textbf{Total 5 a\~nos} & \textbf{\$31,500} & \textbf{\$4,250} & \textbf{\$27,250} \\
\bottomrule
\end{tabular}
\end{table}

\begin{tcolorbox}[colback=secondary!20, colframe=secondary, title=\textbf{Ahorro Proyectado}]
\begin{center}
\textbf{\Huge \$27,250 USD}\\[0.3cm]
\large en 5 a\~nos de operaci\'on
\end{center}
\end{tcolorbox}

%========================================
\section{Rendimiento Esperado}
%========================================

\subsection{Velocidad de Inferencia}

\begin{table}[h]
\centering
\caption{Rendimiento por Configuraci\'on de Hardware}
\begin{tabular}{@{}lccc@{}}
\toprule
\textbf{Modelo} & \textbf{Hardware} & \textbf{Tokens/s} & \textbf{Tiempo respuesta} \\
\midrule
Llama 3.1 8B & 1x RTX 4090 & 80 tok/s & 3-5 segundos \\
Llama 3.1 70B & 2x A100 80GB & 40 tok/s & 8-15 segundos \\
Qwen 2.5 72B & 4x A100 80GB & 35 tok/s & 10-20 segundos \\
\bottomrule
\end{tabular}
\end{table}

\subsection{Capacidad de Usuarios Concurrentes}

\begin{table}[h]
\centering
\caption{Usuarios Simult\'aneos Soportados}
\begin{tabular}{@{}lcc@{}}
\toprule
\textbf{Configuraci\'on} & \textbf{Usuarios simult\'aneos} & \textbf{Latencia p95} \\
\midrule
Ollama + 1 GPU & 3-5 & 15 segundos \\
vLLM + 2 GPU & 10-15 & 12 segundos \\
vLLM + 4 GPU & 25-40 & 10 segundos \\
\bottomrule
\end{tabular}
\end{table}

%========================================
\section{Requisitos de Hardware}
%========================================

\subsection{Configuraci\'on M\'inima}
\begin{itemize}
    \item GPU: 1x NVIDIA con 10+ GB VRAM (RTX 3080, A10)
    \item RAM: 32 GB
    \item CPU: 8 cores
    \item Almacenamiento: 50 GB SSD
    \item Modelo recomendado: Llama 3.1 8B
\end{itemize}

\subsection{Configuraci\'on Recomendada}
\begin{itemize}
    \item GPU: 2x NVIDIA A100 80GB o 4x RTX 4090
    \item RAM: 128 GB
    \item CPU: 32 cores
    \item Almacenamiento: 200 GB NVMe SSD
    \item Modelo recomendado: Llama 3.1 70B
\end{itemize}

\subsection{Configuraci\'on \'Optima}
\begin{itemize}
    \item GPU: 4x NVIDIA A100 80GB o 8x RTX 4090
    \item RAM: 256 GB
    \item CPU: 64 cores
    \item Almacenamiento: 500 GB NVMe SSD
    \item Modelo recomendado: Qwen 2.5 72B / DeepSeek-V2
\end{itemize}

%========================================
\section{Seguridad y Autenticaci\'on}
%========================================

\subsection{Medidas de Seguridad Propuestas}

\begin{enumerate}
    \item \textbf{SSL/TLS}: Todas las comunicaciones cifradas
    \item \textbf{Autenticaci\'on LDAP}: Integraci\'on con directorio universitario
    \item \textbf{Rate Limiting}: M\'aximo 10 consultas/minuto por usuario
    \item \textbf{Firewall}: Solo puerto 443 expuesto al exterior
    \item \textbf{Logs}: Registro de todas las consultas para auditor\'ia
    \item \textbf{Backups}: Respaldo autom\'atico del corpus y configuraciones
\end{enumerate}

\subsection{Privacidad de Datos}

\begin{tcolorbox}[colback=secondary!10, colframe=secondary]
\textbf{Ventaja del despliegue local:} Todas las consultas de los estudiantes permanecen dentro de la infraestructura universitaria. No se env\'ian datos a servidores externos.
\end{tcolorbox}

%========================================
\section{Plan de Implementaci\'on}
%========================================

\subsection{Fases del Proyecto}

\begin{enumerate}
    \item \textbf{Fase 1 - Preparaci\'on}
    \begin{itemize}
        \item Verificar hardware disponible
        \item Instalar sistema operativo y drivers NVIDIA
        \item Configurar red y accesos
    \end{itemize}

    \item \textbf{Fase 2 - Instalaci\'on}
    \begin{itemize}
        \item Instalar Ollama o vLLM
        \item Descargar modelo Llama 3.1
        \item Desplegar ElectroAI
    \end{itemize}

    \item \textbf{Fase 3 - Configuraci\'on}
    \begin{itemize}
        \item Configurar Nginx con SSL
        \item Integrar autenticaci\'on LDAP
        \item Configurar monitoreo
    \end{itemize}

    \item \textbf{Fase 4 - Pruebas}
    \begin{itemize}
        \item Pruebas de carga
        \item Pruebas de seguridad
        \item Pruebas con usuarios piloto
    \end{itemize}

    \item \textbf{Fase 5 - Producci\'on}
    \begin{itemize}
        \item Despliegue gradual
        \item Capacitaci\'on a estudiantes
        \item Monitoreo continuo
    \end{itemize}
\end{enumerate}

%========================================
\section{Conclusiones y Recomendaciones}
%========================================

\subsection{Conclusiones}

\begin{enumerate}
    \item \textbf{Factibilidad}: El despliegue local es \textbf{completamente factible} con el hardware disponible.

    \item \textbf{Costo-beneficio}: El ahorro proyectado de \textbf{\$27,250 USD en 5 a\~nos} justifica ampliamente la inversi\'on de tiempo en implementaci\'on.

    \item \textbf{Calidad}: Los modelos open source como Llama 3.1 70B ofrecen calidad \textbf{comparable a Claude} para tareas educativas.

    \item \textbf{Control}: La universidad mantiene \textbf{control total} sobre el servicio, datos y disponibilidad.

    \item \textbf{Escalabilidad}: El sistema puede escalar agregando m\'as GPUs sin costos adicionales de licencia.
\end{enumerate}

\subsection{Recomendaci\'on Final}

\begin{tcolorbox}[colback=primary!20, colframe=primary, title=\textbf{Recomendaci\'on}]
\textbf{Se recomienda proceder con el despliegue local} utilizando:
\begin{itemize}
    \item Modelo: \textbf{Llama 3.1 70B} (mejor balance calidad/recursos)
    \item Servidor: \textbf{Ollama} (facilidad de uso) o \textbf{vLLM} (mejor rendimiento)
    \item Autenticaci\'on: \textbf{LDAP universitario}
\end{itemize}

Esta configuraci\'on permitir\'a atender a \textbf{500+ estudiantes} con un costo operativo de aproximadamente \textbf{\$70 USD/mes}.
\end{tcolorbox}

\vspace{1cm}

\begin{center}
\textcolor{primary}{\rule{0.5\linewidth}{1pt}}\\[0.5cm]
\textit{Documento generado por ElectroAI}\\
\textit{Universidad del B\'io-B\'io}\\
\today
\end{center}

\end{document}
